% ch09: 原文 GLDraft_D.tex 行 1974–2284；§6
% 只译散文（标题、正文、图题、脚注）；公式、\label、\ref、\eqref、\cite、图文件名保持不变。

\section{实丛}\label{realb}

\subsection{设置}\label{setup}

到目前为止，我们一直把 $\M_H(G,C)$ 作为实辛流形来研究其量子化。另一种做法是：把 $\M_H(G,C)$ 视为带有复结构 $I$ 与全纯辛结构 $\Omega_I=\omega_J+\i \omega_K$ 的复辛流形，转而量子化一个实辛子流形 $\M_H^\R \subset \M_H$。为此，如第 \ref{background} 节所解释的，我们要找 $\M_H(G,C)$ 的一个反全纯对合（antiholomorphic involution）$\tau$，使之满足 $\tau^*\Omega_I=\bar\Omega_I$。在我们的应用中，$\M_H^\R$ 将是一个余切丛，从而使希尔伯特空间与 Hecke 算符的作用都有简单的描述；它还将是一个带有 Hitchin 纤维化实版本的实可积系统（real integrable system）。

可能的反全纯对合的分类，无论所考虑的是 \cite{Hurtu,Hurtu2} 中那样的全纯 $G_\C$ 丛，还是 \cite{BaSc,BaSc2,BG} 中那样的希格斯丛，都是一样的。这些文献所提供的细节远多于我们此处将要解释的。我们还要指出，在研究对偶时将会用到的三维流形 $U_\tau=(\Sigma\times \hat I)/\{1,\tau\}$，是在 \cite{BaSc} 第 11 节中引入的。该文还基于 $\M_H^\R(G,C)$ 作为实可积系统的结构，提出了一个对偶方案。

$\M_H$ 的适当对合可以从黎曼曲面 $C$ 的反全纯对合 $\tau$ 出发来构造；对适当选择的 $C$，这样的对合是存在的。反全纯映射反转 $C$ 的定向；反过来，若 $\tau$ 是可定向二维流形 $C$ 上的反定向光滑对合，则可以选取 $C$ 上的复结构，使得 $\tau$ 反全纯地作用。\footnote{在 $C$ 上任取一个黎曼度量 $\mathrm g$。则 $\mathrm g'=(\mathrm g+\tau^*(\mathrm g))/2$ 是一个 $\tau$ 不变度量，而在由 $\mathrm g'$ 决定的复结构中，$\tau$ 反全纯地作用。}

拓扑上，$\tau$ 的可能选择可以如下分类。若 $C$ 的亏格为 $g$，则其欧拉示性数为 $2-2g$，而商 $C'=C/\{1,\tau\}$ 的欧拉示性数为 $1-g$。$C'$ 可以是欧拉示性数为 $1-g$ 的任意（可能不可定向的、可能带边界的）二维流形。$C'$ 的边界来自 $\tau:C\to C$ 的不动点。这些不动点组成若干个圆周；圆周的数目可以是 $0$ 到 $g+1$ 之间的任意整数。例如，若 $g=0$，$C'$ 可以是圆盘或 $\Bbb{RP}^2$；若 $g=1$，$C'$ 可以是圆柱面、Möbius 带或 Klein 瓶。

固定 $\tau$ 之后，我们接着选择 $\tau$ 提升到光滑 $G$-丛 $E\to C$ 上作用的拓扑型。我们选择保持 $\tau^2=1$ 的提升。若 $s$ 是 $\tau:C\to C$ 的不动点，则 $\tau$ 通过 $G$ 的自同构 $x_s$ 作用在 $E$ 于 $s$ 处的纤维 $E_s$ 上。这里 $x_s$ 可以是内自同构，即由 $G$ 中元素 $g_s$ 所作的共轭；但更一般地，若 $G$ 有非平凡外自同构，$x_s$ 也可以是外自同构。$x_s$ 的共轭类在每个不动圆 $S$ 上是常值的，我们把它记作 $x_S$。条件 $\tau^2=1$ 对 $x_S^2$ 施加了一个条件。

由于希格斯丛由伴随值的场 $(A,\phi)$ 定义，仅仅为了定义希格斯丛模空间 $\M_H(G,C)$ 的一个对合，我们只要求 $x_S^2=1$ 在 $G$ 的伴随形式中成立；作为 $G$ 的元素，$x_S^2$ 可能是一个不等于 1 的中心元素。应当对 $x_S^2$ 施加的确切条件，取决于第 \ref{toposubt} 节所回顾的那些拓扑微妙之处；我们将在第 \ref{nonfree} 节回到这一点。

若 $\tau$ 自由作用且 $G$ 不单连通，则在把 $\tau$ 提升为对 $E$ 的作用时会出现其他问题。例如，对 $G=\SO(3)$，只有当 $\int_C w_2(E)=0$ 时，$\tau$ 才存在提升作用在 $E$ 上。在这种情形，有两个拓扑上不等价的提升，由 $\int_{C'}w_2(E')$ 参数化，其中 $E'\to C'$ 是 $E\to C$ 对 $\{1,\tau\}$ 的商。

带有反全纯对合 $\tau$ 的黎曼曲面 $C$ 可以看作定义在 $\R$ 上的代数曲线。$C'$ 的边界点（若有）对应于 $C$ 在 $\R$ 上的实点。一旦 $\tau$ 被提升为作用在光滑丛 $E$ 上，它就也作用在 $\M(G,C)$（即 $E$ 上平坦联络的模空间）以及相应的希格斯丛模空间 $\M_H(G,C)$ 上。无论 $\tau$ 在 $C$ 上的作用是否有不动点，$\tau$ 在\footnote{若 $A$ 是 $E$ 上任意联络，则 $\frac{1}{2}(A+\tau^*(A))$ 是 $\tau$ 不变联络。一般而言，该联络的 $(0,1)$ 部分定义一个稳定 $G_\C$-丛，它对应于 $\M(G,C)$ 中的一个 $\tau$ 不变点。同一个联络配上零希格斯场，则定义 $\M_H(G,C)$ 中的一个 $\tau$ 不变点。}$\M(G,C)$ 上以及在 $\M_H(G,C)$ 上的作用总有不动点。

$\tau$ 在 $\M(G,C)$ 或 $\M_H(G,C)$ 上不动点集的每个连通分支都是中间维数（middle-dimensional）的。这是复流形的反全纯对合的一般性质。\zy{译者按（无外部依据）：在 $\tau$ 的不动点处，切映射 $\d\tau$ 是实切空间上的对合，并且反全纯（即与复结构 $I$ 反交换）；因此它的 $+1$ 本征空间恰有复维数那么多维，不动点集的每个连通分支都是中间维数的——其实维数等于复维数，即整个实维数的一半。}

我们将用 $\M^\R(G,C)$ 或 $\M_H^\R(G,C)$ 分别表示 $\tau$ 作用在 $\M(G,C)$ 或 $\M_H(G,C)$ 上不动点集的一个连通分支。这些连通分支一般而言由我们尚未引入的附加数据来分类。特别地，把 $\M(G,C)$ 或 $\M_H(G,C)$ 视为 $G$ 值或 $G_\C$ 值平坦联络的模空间时，平坦丛绕 $S$ 的单值 $h_S$ 在 $x_S$ 下不变，因此它落在 $G$ 或 $G_\C$ 的 $x_S$ 不变子群 $H_S$ 或 $H_{S,\C}$ 中。若 $G$ 单连通，则 $H_S$ 与 $H_{S,\C}$ 连通；但一般情形并非如此（例如，若 $G=\SO(3)$ 且 $g_S=\mathrm{diag}(1,-1,-1)$，则 $H_S$ 包含单位元连通分支以及另一个包含元素 $\mathrm{diag}(-1,-1,1)$ 的连通分支）。所以一般而言，要指定不动点集的一个连通分支，需要对每个 $S$ 作额外的选择。

与每个不动圆 $S$ 相伴的是复李群 $G_\C$ 的一个实形式（real form）$G_{\R,S}$。把 $g\to \bar g$ 记作 $G_\C$ 的、保持紧形式 $G$ 不变的反全纯对合，则 $G_{\R,S}$ 由条件
\be\label{ralform} g=x_S(\bar g). \ee
定义。当 $x_S$ 是由 $g_S\in G$ 所作的共轭时，该条件变为
\be\label{realform} g=g_S \bar g g_S^{-1}.\ee

带有辛结构 $\omega_J$ 的 $\M_H^\R(G,C)$，就是我们想要量子化的实辛流形。$\M_H^\R(G,C)$ 的定义依赖于若干我们未在记号中标明的选择；但无论作何选择，$\M_H^\R(G,C)$ 都以 $\M_H(G,C)$ 为其复化，而这个复化具有适合对 $\M_H^\R(G,C)$ 作膜量子化的性质。

$\M_H^\R(G,C)$ 有一个实极化（real polarization），定义如下：极化的一个叶由 $\tau$ 不变的希格斯对 $(A,\phi)$ 组成，其中 $A$ 固定而 $\phi$ 变化。同样的条件若不要求 $\tau$ 不变性，就定义 $\M^\R_H(G,C)$ 的复化 $\M_H(G,C)$ 的一个全纯极化（holomorphic polarization）。因此，$\M^\R(G,C)$ 的膜量子化等价于用这个实极化所作的量子化。

具体地说，在这个实极化中，$\M_H^\R(G,C)$ 可以用余切丛 $T^*\M^\R(G,C)$ 来近似，理由与完整的希格斯丛模空间有类似陈述的理由相同：若 $(A,\phi)$ 是表示 $\M_H^\R(G,C)$ 中一点的 $\tau$ 不变希格斯丛，则一般而言 $\phi$ 表示 $\M^\R(G,C)$ 在对应于 $A$ 的那一点处的余切向量。

因此，在对 $\M^\R_H(G,C)$ 作量子化时出现的希尔伯特空间 $\H_\tau$，就是 $\M^\R(G,C)$ 上 $\lmark^2$ 半密度的空间——这是余切丛几何量子化的通常答案。

我们还可以把 Hitchin 纤维化限制到 $\tau$ 不变轨迹上。$\tau$ 作用在 Hitchin 纤维化的底空间 ${\cB}$ 上，其不动点集为 ${\cB}^\R$。${\cB}^\R$ 上的函数彼此泊松交换（Poisson-commuting，相对于 $\M_H^\R(G,C)$ 的实辛结构 $\omega_J$）。因此 $\M_H^\R(G,C)$ 是一个实可积系统。映射 $\M_H^\R(G,C)\to {\cB}^\R$ 的一般纤维是一个实环面。Baraglia 和 Schaposnik \cite{BaSc} 提议，通过在实可积系统的纤维上作 $T$-对偶，在 $\M_H^\R(G,C)$ 与把 $G$ 换成 $G^\vee$ 的类似模空间之间定义镜像对称。一般而言，人们会预期这样的定义至少在渐近意义下、即在远离实可积系统判别式轨迹之处给出可靠的结果。至少在 $\tau$ 自由作用时，人们可以更精确一些；我们下面就来讨论。

\subsection{四维图景与对偶}\label{fourdpic}

为了从这一构造中得到有趣的结果，人们想要应用对偶；为此，一如往常，四维图景是有帮助的。

对于四维图景，我们从通常的四维流形 $M=\Sigma\times C$ 出发，其中 $\Sigma=\R\times I$ 是一条带。这里 $\R$ 标记“时间”，$I$ 是一个区间，比如说 $0\leq w\leq 1$。然后我们通过对 $w=1$ 处的右边界施加一个等价关系来修改 $M$：我们规定，对 $p\in C$，点 $(t,1,p)$ 与 $(t,1,\tau(p))$ 等价。（除 $w=1$ 外不施加任何等价关系。）施加这个等价关系相当于要求 $w=1$ 处的数据是 $\tau$ 不变的（相差一个规范变换）。因此，这是在四维中实现第 \ref{setup} 节二维图景中右边界处膜 $\B_\R$ 的一种方式。我们把 $M$ 对这个等价关系的商记作 $M_\tau$。方便的做法是把时间分离出来，定义 $M_\tau=\R\times U_\tau$。

$U_\tau$ 还有另一种构造，即作为由 $\tau$ 生成的 $\Z_2$ 群的商。为此，我们从加倍的区间 $\h I: 0\leq w\leq 2$ 出发。然后把 $\h U= \h I\times C$ 按对称 $(w,p)\to (2-w,\tau(p))$ 作商。基本域是区域 $0\leq w\leq 1$，所以商与前面相同。

最简单的情形是 $\tau$ 在 $C$ 上自由作用。此时 $U_\tau$ 实际上是一个可定向流形，其边界是 $C$ 的一份拷贝，位于 $w=0$ 处。$U_\tau$ 是 $C'=C/\{1,\tau\}$ 上一个实线丛的全空间。$C'$ 本身与这个实线丛都不可定向，但全空间 $U_\tau$ 可定向。

与 $\M_H^\R(G,C)$ 的量子化相联系的四维图景，就是 $M_\tau=\R\times U_\tau$ 上的 $A$-模型，在 $w=0$ 处取 $\B_\cc$ 边界条件。在 $w=1$ 处不需要显式的边界条件，因为那里没有边界。二维描述中支撑在 $\M_H^\R(G,C)$ 上的拉格朗日膜 $\B_\R$，已被吸收进 $M_\tau$ 的几何之中。

因此，我们可以径直描述一个对偶的图景。对偶就是同一个四维流形 $M_\tau$ 上规范群为 $G^\vee$ 的 $B$-模型，但此时在 $w=0$ 处取 $\B_\op$ 边界条件。用二维语言说，与 $\B_\R$ 对偶的膜 $\B^\vee_\R$ 支撑在 $\M_H^\R(G^\vee,C)$ 上。特别地，它支撑在 Hitchin 纤维化底空间上同一个轨迹 ${\cB}_\R$ 上。这与 Baraglia 和 Schaposnik \cite{BaSc} 的对偶方案相符；该方案似乎至少在 $\tau$ 自由作用于 $C$ 时成立。

一如往常，在 $U_\tau$ 上量子化时，描述 $B$-模型的物理态以及 Hitchin 哈密顿量的本征值相对容易。第一步是把 $B$-模型局域化到复值平坦联络上，这里指的是 $U_\tau$ 上满足 oper 边界条件的 $G^\vee_\C$ 值平坦联络。换一种说法，局域化到平坦 $G_\C$ 丛 $E^\vee\to U_\tau$ 上，使其在 $C=\partial U_\tau$ 上的限制是一个全纯 oper。令 $\Upsilon_\tau$ 为这类丛的同构类的集合。为了希尔伯特空间内积的正定性，我们预期 $\Upsilon_\tau$ 是非退化点的离散集（见附录 \ref{bmodel}），类似于在完整希格斯丛模空间量子化中遇到的集合 $\Upsilon=L_\op\cap L_\bop$。假定如此，希尔伯特空间 $\H_\tau$ 具有由 $u\in \Upsilon_\tau$ 标记的基 $\psi_u$。这些 $\psi_u$ 是 Hitchin 哈密顿量的本征函数。Hitchin 哈密顿量 $H_{\P,\alpha}$ 在 $\psi_u$ 上的本征值，就是相应函数 $f_{\P,\alpha}:L_\op\to \C$ 在 $L_\op$ 中对应于 $u$ 所关联平坦丛的边界值的那一点处的取值。

若 $C=\partial U_\tau$ 上的一个 oper 丛作为平坦丛延拓到 $U_\tau$ 上，则特别地，这意味着反全纯对合 $\tau:C\to C$ 提升为 $E^\vee$ 上的作用，从而 $E^\vee$ 既是一个反全纯 oper，又是一个全纯 oper。\zy{译者按（无外部依据，据文义校勘）：原文此处作 “antihomomorphic involution”，当为 “antiholomorphic involution”（反全纯对合）之误；译文按反全纯对合处理，特此说明。} 所以集合 $\Upsilon_\tau$ 有一个自然映射到 $\Upsilon=L_\op\cap L_\bop$。然而，一般而言从 $\Upsilon_\tau$ 到 $\Upsilon$ 的映射不是嵌入，因为 $U_\tau$ 边界上作为全纯 oper 的 $\tau$ 不变平坦丛，作为 $U_\tau$ 内部上的平坦丛可能有不止一个延拓。为分析这一情形，注意 $U_\tau$ 上的平坦 $G^\vee_\C$ 丛对应于同态 $\varrho:\pi_1(U_\tau)\to G^\vee_\C$（模去共轭）。由于 $U_\tau$ 可缩到 $C'=C/\{1,\tau\}$，我们同样可以考虑 $\varrho:\pi_1(C')\to G^\vee_\C$。群 $\pi_1(C')$ 有一个指数 2 的子群 $\pi_{1,+}(C')$，由 $C'$ 中保定向的环路组成。这样的环路可以形变到 $U_\tau$ 的边界，所以一旦在 $U_\tau$ 的边界上给定平坦 $G^\vee_\C$ 丛，相应同态 $\varrho$ 在 $\pi_{1,+}(C')$ 上的限制就唯一确定。然而，$\varrho$ 到 $\pi_1(C')$ 其余部分的延拓有一定的自由。$\pi_1(C')$ 由指数 2 子群 $\pi_{1,+}(C')$ 与任一反定向元素 $\sigma\in \pi_1(C')$ 生成。为完成对 $\varrho$ 的描述，我们需要指定 $\varrho(\sigma)$。$\varrho$ 应当是同态，所以我们要求 $\varrho(\sigma)^2=\varrho(\sigma^2)$；由于 $\sigma^2\in \pi_{1,+}(C')$，$\varrho(\sigma^2)$ 由边界数据决定。若 $\varrho(\sigma^2)$ 是 $G^\vee_\C$ 的正则元素，则 $\varrho(\sigma)$ 只有有限多个选择，但 $\varrho(\sigma)$ 并不唯一确定。特别地，我们总可以作变换 $\varrho(\sigma)\to \veps \varrho(\sigma)$，其中 $\veps$ 是 $G^\vee$ 的中心的一个 2 阶元。于是，$G^\vee$ 的中心中由 2 阶元组成的子群 $\zZ_2(G^\vee)$ 自由地作用在 $\Upsilon_\tau$ 上。（若 $G^\vee$ 单，则 $\zZ_2(G^\vee)$ 为 1、$\Z_2$ 或 $\Z_2\times \Z_2$，视 $G^\vee$ 而定。）

$\zZ_2(G^\vee)$ 作用的 $A$-模型对偶如下。拓扑上，当 $G$ 单时，三维流形 $U_\tau$ 上的 $G$-丛由 $H^2(U_\tau,\pi_1(G))$ 分类。在目前情形，$U_\tau$ 可缩到 $C'=C/\tau$，故 $H^2(U_\tau,\pi_1(G))=H^2(C',\pi_1(G))$。由于 $C'$ 不可定向，有 $H^2(C',\pi_1(G))=\pi_{1;2}(G)$，其中 $\pi_{1;2}(G)$ 是 $\pi_1(G)$ 中由 2 阶元组成的子群。\zy{译者按：$\tau$ 自由作用、$C'$ 闭不可定向时，局部系数 Poincaré 对偶给 $H^2(C';\pi_1(G))\cong\pi_1(G)/2\pi_1(G)$（Hatcher 第 3.H 节）。原文进一步将其等同于 2-挠子群 $\pi_{1;2}(G)$；这一等同一般只在 $\pi_1(G)$ 为 2-群时成立（例如 $G=E_6$、$\pi_1(G)=\Z/3$ 时 $\pi_1(G)/2\pi_1(G)=\Z/3$，而 $\pi_{1;2}(G)=0$）。} 于是，在 $A$-模型一侧，希尔伯特空间有一个由 $\pi_{1;2}(G)$ 给出的分次。由于 $\pi_{1;2}(G)=\zZ_2(G^\vee)$，这与 $B$-模型中的 $\zZ_2(G^\vee)$ 作用相匹配。

最后让我们讨论这个问题中的自然线算符。如第 \ref{wtw} 节那样，原则上我们可以考虑支撑在任意一维流形 $\gamma\subset M_\tau$ 上的 Wilson 与 't Hooft 线算符。但一个自然的特殊情形如下。取一点 $p\in C$。在 $\h I\times C$ 中，有一条自然的路径 $\h\gamma_p$，连接左、右边界上的点 $0\times p$ 与 $2\times p$：我们直接取 $\h\gamma_p = \h I\times p$。对由 $\tau$ 生成的群作商后，$\h\gamma_p$ 下降到 $U_\tau$ 中的一条路径 $\gamma_p$，连接边界点 $0\times p$ 与 $0\times \tau(p)$。在纯粹的二维描述中，人们需要与实膜 ${\B}_\R$ 相关联的端点或角数据；但在四维描述中，这并不需要，因为 $\gamma$ 在 $w=1$ 处没有端点。

如第 \ref{wilop} 节，我们可以选取 $G^\vee$ 的任意表示 $R$，并考虑丛 $E^\vee_R=E^\vee\times_{G^\vee_\C}R$ 在曲线 $\gamma_p$ 上的和乐 $W_{R,\gamma}$。然而，要使这个和乐成为一个量子算符 $W_{R,p}$，我们需要像以前一样提供端点数据。例如，若我们在 $0\times p$ 处选取典范端点 $s_{R,n}$，在 $0\times \tau(p)$ 处选取 $s_{\bar R, m}$，则 Wilson 算符将取为输运到 $\tau(p)$ 的 $s_{R,n}$ 与 $s_{\bar R, m}$ 之间内积的形式。如第 \ref{wilop} 节，这样构造的 Wilson 算符 $W_{R,p}$ 在态基 $\psi_u$ 上是对角的，其本征值由在对应于 $u$ 的平坦丛中对具体的典范截面求内积给出。

对偶预言：与同一表示 $R$ 相关联、支撑在同一条曲线 $\gamma_p$ 上、并带有 $S$-对偶端点的对偶 't Hooft 算符，具有相同的本征值。对偶 't Hooft 算符与 $\M_H^\R(G,C)$ 到自身的实 Hecke 对应（real Hecke correspondence）$Z^\R_{R,p}$ 相关联；$Z^\R_{R,p}$ 就是 $\M_H(G,C)$ 到自身的通常 Hecke 对应 $Z_{R,p;\bar R,\tau(p)}$ 的 $\tau$ 不变轨迹，而该通常对应是对一对 $\tau$ 共轭点、以共轭（或对偶）表示对 $R,\bar R$ 为标签而定义的。如第 \ref{lineq} 节，由于余切丛的标度对称性，与 Hecke 对应相关联的量子算符可以由半经典公式定义。除第 \ref{lineq} 节所用到的要素之外，还需要一个进一步的事实。设 $X$ 是带有典范丛 $K_X$ 的复流形；令 $\tau:X\to X$ 是反全纯对合，其不动点集为 $X^\R$，并假设 $X^\R$ 可定向。\zy{译者按（无外部依据，据文义校勘）：原文 “and assume thast $X^\R$ is orientable” 中的 “thast” 系 “that” 之误；译文按“假设 $X^\R$ 可定向”处理。} 于是 $K_X^{1/2}\to X$ 的全纯截面在 $X^\R$ 上限制为复值半密度。因此，人们在定义 $\M_H$ 量子化中的 't Hooft 算符时会用到的（连同反全纯端点或角一起的）全纯端点或角，在实轨迹上限制为一个半密度，它定义了 $\M_H^\R$ 量子化中的 't Hooft 算符。在我们的应用中，$Z^\R_{R,p}$ 可定向，因为 $Z_{R,p;\bar R,\tau(p)}$ 是 Calabi–Yau 流形，如第 \ref{lineq} 节所讨论；具有实结构的 Calabi–Yau 流形 $X$ 的 Calabi–Yau 形式可以选为实的，并在实轨迹上限制为定义 $X^\R$ 定向的最高次微分形式。

\subsection{$\tau$ 不自由作用的情形}\label{nonfree}

现在考虑 $\tau$ 在 $C$ 上不自由作用的情形。设 $\tau$ 的作用固定了一个圆周 $S\subset C$。则 $U_\tau=\h U/\Z_2$ 把 $S$ 作为 $\Z_2$ 轨形（orbifold）不动点的轨迹包含在内。$S$ 附近的局部行为形如
\be\label{wofo} S\times \R^2/\{1,\tau\}, \ee
其中 $\tau$ 在 $\R^2$ 上作为 $\pi$ 旋转作用。在四维流形 $M_\tau=\R\times U_\tau$ 中，不动点集是 $\R\times S$。

我们回忆，一般而言 $\tau$ 在 $S$ 上的作用伴随着自同构 $x_S$ 的作用，后者至少在 $G$ 的伴随形式中满足 $x_S^2=1$。在 $G$ 规范场论中，人们也许预期要求在 $G$ 中 $x_S^2=1$，但必须考虑第 \ref{toposubt} 节所回顾的拓扑微妙之处。

为说明这些问题，我们考虑 $G$ 秩为 1、从而是 $\SU(2)$ 或 $\SO(3)$ 的情形。此时 $G$ 没有外自同构，$x_S$ 是由 $G$ 中元素 $g_S$ 所作的共轭。如第 \ref{toposubt} 节所讨论，$\SU(2)$ 规范场论有两个版本。在通常的 $\SU(2)$ 规范场论中，最明显的条件是要求 $g_S^2=1$ 作用在 $\SU(2)$ 丛 $E\to M$ 上。于是仅有的选项是 $g_S=1$ 与 $g_S=-1$。另一方面，在 $\Spin\cdot \SU(2)$ 规范场论中，最明显的条件是要求 $g_S^2=1$ 作用在 $\cS\otimes E$ 上，其中 $\cS$ 是 $M$ 上 $\tau$ 不变的旋量丛（spin bundle，至少局部定义在不动点集附近），$E$ 是 $\SU(2)$ 丛（在 $\cS$ 有定义处有定义）。由于 $\tau$ 在不动点集处作为法平面的 $\pi$ 旋转作用，$\tau^2$ 是 $2\pi$ 旋转，并在 $\cS$ 上作为 $-1$ 作用。这似乎提示我们应要求 $\tau^2=-1$ 作用在 $E$ 上，从而（在共轭意义下）$g_S=\mathrm{diag}(\i,-\i)$。\footnote{这一共轭类的作用是由 D. Baraglia 提示的。}

然而，我们认为也可能可以把这些选择反过来。例如，在 $\SU(2)$ 中，在试图对 $\tau$ 作商之前，我们可以假定沿 $\tau$ 的不动点轨迹有一个单值缺陷（monodromy defect），其单值为 $-1$。于是在作商时我们会想要 $g_S^2=-1$。类似地，在作商之前加入这样的单值缺陷，会为 $\Spin\cdot \SU(2)$ 提示 $g_S^2=+1$。

类似地，$\SO(3)$ 规范场论也有两个版本，区别在于路径积分被积函数中是否含因子 $\Delta=(-1)^{\int_M w_2(M) w_2(E)}$。无论是否含这个因子，我们都想要 $g_S^2=1$，这给出两种可能，即 $g_S=1$ 与 $g_S=\mathrm{diag}(-1,-1,1)$。在存在 $g_S=\mathrm{diag}(-1,-1,1)$ 的余维 2 奇点时，$\Delta$ 在拓扑上没有良定义，所以在此情形看来我们想要 $g_S=1$。若不含因子 $\Delta$，则 $g_S$ 的两种可能都可行。对 $\SO(3)$，由于中心平凡，我们没有加入带中心单值的缺陷这一选项，但可以加入它的对偶，即一个感知限制到不动点集上的规范丛 $E$ 的拓扑的缺陷。对 $\SO(3)$，这个缺陷是路径积分中的一个因子 $(-1)^{\int_W w_2(E)}$，其中 $W$ 是不动点集。

一般而言，找到这类奇点的对偶是一个微妙的问题。之所以微妙，部分原因在于：人们{\it 不能}假定，由 $\tau$ 不变性条件所定义的 $\R^2/\Z_2$ 轨形奇点，其对偶是由对偶的 $\tau$ 不变性条件所定义的另一个 $\R^2/\Z_2$ 轨形奇点。一般而言，人们只知道 $\R^2/\Z_2$ 轨形奇点的对偶是一个余维 2 缺陷，它保持与 $\R^2/\Z_2$ 轨形奇点相同的超对称。在一个有些类似的刚性问题，即刚性面缺陷（rigid surface defects）的问题中，要对对偶的作用获得一般性的理解被证明是困难的 \cite{GWr}。

\begin{yizhu}[译者按（无外部依据）：轨形奇点与它的对偶为何“微妙”]
$\R^2/\Z_2$ 中的 $\Z_2$ 由 $\pi$ 旋转生成；商空间作为拓扑空间仍是 $\R^2$，奇性只体现在轨形结构上：原点是轨形点，其“切空间”要按 $\Z_2$ 的表示来理解。$\tau$ 在 $C$ 上有不动圆 $S$ 时，$U_\tau$ 沿 $S$ 正出现这种轨形奇点（式 \eqref{wofo}）。本文在 $A$-模型一侧考虑用光滑流形 $\t U_\tau$ 代替 $U_\tau$，但原文是以“如果这一步是正确的”这样的条件句提出的。在 $B$-模型一侧，对偶的对象一般不再是同类型的 $\R^2/\Z_2$ 轨形奇点，而只是一个保持相同超对称的余维 2 缺陷，因而不能直接由“$\tau$ 不变性”的条件读出。作者因此称这是一个“微妙”的问题，只给出一个原则上分析一般情形的策略；这与第 \ref{fourdpic} 节中关于 Baraglia--Schaposnik 提议“至少在 $\tau$ 自由作用时成立”的限定是一致的。
\end{yizhu}

如果所有 $S$ 的 $g_S$ 都是中心的，人们会忍不住根据以下理由断言一个简单的答案。设 $g_S=1$。那么奇性只存在于几何中，而不在规范场中。$\R^2/\Z_2$ 轨形奇点在拓扑上根本不是奇性，因为 $\R^2/\Z_2$ 在拓扑上等价于 $\R^2$。所以若对所有不动圆都有 $g_S=1$，则 $U_\tau$ 在拓扑上实际上是一个流形，人们可以把轨形奇点“磨圆”，赋予它一个光滑的几何。我们把磨圆后的版本称为 $\t U_\tau$。如果在 $A$-模型中用光滑流形 $\t U_\tau$ 代替轨形是正确的，那么对偶就是同一个流形上的 $B$-模型。这一推理对更一般的情形有一个潜在的类比：$g_S$ 是中心的但不等于单位元。我们仍可以把缺陷磨圆以得到光滑流形 $\t U_\tau$，但此时 $\t U_\tau$ 含有支撑在轨形轨迹 $W$ 上的带中心单值的缺陷。如前面所说，这个缺陷的对偶是一个感知丛的拓扑的缺陷，例如在规范群为 $\SO(3)$ 的情形由因子 $(-1)^{\int_W w_2(E)}$ 定义的缺陷。

下面原则上是一个分析一般情形的策略。用二维环面把不动点轨迹围起来。于是轨形缺陷定义了由四维超杨–米尔斯理论经 $T^2$ 紧化所得的二维理论的一个边界条件。由于有大的未破缺规范对称性，这并不完全是一个 $\sigma$-模型。如果我们能识别该边界条件的镜像，它就会给出对偶一侧限制在二维环面上的 $G^\vee_\C$ 平坦联络的边界条件。这将足以刻画 $U_\tau$ 边界处的 oper 平坦联络如何延拓到内部，从而确定对偶 $B$-模型一侧的谱。

\subsection{实 Hecke 算符}\label{rho}

$\tau$ 在 $C$ 上不自由作用的情形最重要的新特点，也许是存在支撑在 $C$ 中一个实点上的线算符，这与第 \ref{fourdpic} 节所考虑的支撑在一对 $\tau$ 共轭点上的线算符相对。

首先我们描述支撑在实点 $p\in C$ 上的 't Hooft 算符的规范场论图景。我们在四维流形 $M=\R\times \h I\times C$ 上工作，其中 $\R$ 由 $t$ 参数化，$\h I$ 由 $w$ 参数化。局部图景即已足够，所以我们取 $C$ 就是复 $z$-平面 $\C$，并考虑由 $(t,w,z)\to (-t,w,\bar z)$ 作用的对合 $\tau$。仅作用在 $\C$ 上时，$\tau$ 有由 ${\mathrm{Im}\, z}=0$ 定义的不动线 $S$；我们让 $\tau$ 伴随一个满足 $x_S^2=1$ 的自同构 $x_S$（或第 \ref{nonfree} 节所讨论的稍更一般的条件）。支撑在 $z=t=0$ 的 't Hooft 算符由 $G$ 规范场论的 Dirac 单极子解描述。该解的结构群约化为极大环面 $T\subset G$，并可由其曲率刻画：
\be\label{curv} F=\frac{\m}{2}\star_3 \d \frac{1}{(t^2+|z|^2)^{1/2}}, \ee
其中 $\star_3$ 是 $\R\times\C$ 上度量 $\d t^2+|\d z|^2$ 的 Hodge 星，$\m$ 是 $T$ 的李代数 $\frak t$ 中的一个常元素。要使具有此曲率的联络存在，$\m$ 必须是整余权（integral coweight，即余权格中的元素），对偶于 $G^\vee$ 的表示 $R$（它与第 \ref{affgr} 节中称为 $\m$ 的对象一致）。现在我们要问：这个解在 $\tau$ 连同自同构 $x_S$ 的作用下是否不变？由于 $\star_3$ 在 $\tau$ 下为奇，一个必要且充分的条件是 $\m$ 在 $x_S$ 下为奇：
\be\label{oddurv} x_S( \m )=-\m. \ee
当且仅当可以在 $\m$ 的共轭类中选取 $\m$ 使其在 $x_S$ 下为奇时，这样构造的解才是 $\tau$ 不变的，并下降到 $M_\tau=M/\{1,\tau\}$ 上的一个解。它描述一个实 't Hooft 算符，支撑在 $C$ 中的一个实点上，并与 $G^\vee$ 的表示 $R$ 相关联。

利用这个模型解，我们可以定义“实”Hecke 修正的空间，它可以由实 't Hooft 算符实现。存在 $G_{\C}({\cal K})$ 与 $G_{\C}({\cal O})$ 的实版本 $G_{\R,S}({\cal K})$ 与 $G_{\R,S}({\cal O})$，它们由沿 $S$ 取值于 $G_{\R,S}$ 的规范变换组成。于是我们可以定义仿射 Grassmann 流形的实版本
\begin{equation}
\Gr_{G_{\R,S}} = G_{\R,S}({\cal K})/G_{\R,S}({\cal O})
\end{equation}
以及类型为 $\m$ 的实 Hecke 修正（real Hecke modification）的轨道 $\left[G_{\R,S}({\cal O}) z^\m\right]$。

知道给定 $x_S$ 下哪些实 't Hooft 算符是可能的，会对具有给定 $x_S$ 的 $\Z_2$ 轨形奇点的对偶施加非常强的约束。

几何上，实 't Hooft 算符从 $U_\tau$ 的边界延伸到内部的一个不动点 $p$。对实 't Hooft 算符的完整分析（我们在此不打算做）将包括讨论 't Hooft 算符在不动点上的可能端点，以及推导相应的积分算符。

在 $S$-对偶一侧，我们将有一个从边界延伸到 $p$ 的 Wilson 算符。Wilson 线的一个特定端点会给出 $p$ 附近规范丛的平坦截面空间中的一个向量 $v_{\bar R,m}$。若不知道轨形奇点的 $S$-对偶，我们就无法更精确地刻画这个向量；特别地，$B$-模型中的平坦丛可能不会延拓到不动点集上。Wilson 算符的期望值将取配对 $(v_{\bar R,m}, s_{R,n})$ 的形式，从而选出 oper 微分方程的一个特定解。

更一般地，可以考虑带有边界 $C$ 的任意可定向三维流形 $U$，并研究 $\R\times U$ 上的对偶 $A$-模型与 $B$-模型。抛开技术困难（$U$ 上平坦丛的模空间可能非常奇异），人们可以希望定义物理态的空间 $\H_U$，在 $A$-模型一侧带有 Hitchin 哈密顿量的作用，并在 $B$-模型一侧用经典数据预言它们的本征值。人们也可以研究线算符，不过一般而言不会有与本文所考虑的那些算符密切类似的算符。态空间 $\H_U$ 将带有埃尔米特内积（hermitian inner product），但一般而言并不清楚这个内积是否正定，因为它不再是通过量子化余切丛得到的。
